\cotefile{64}
\scanpage{40}
% archive sheet number: 39
\(\underline{\mathrm I}\unsure{c|\subset}\)

\uline{NB} \(X\mid\underline J\simeq\Delta\) est la configuration octaédrale\ldots

\diagram{\resizebox{\linewidth}{!}{\begin{tikzcd}[ampersand replacement=\&,column sep=small,row sep=small]
\begin{gathered}(X,\underline I)\ \text{\deleted{\unsure{étale|début}}}\ 
\underline I\ \text{\unsure{rev.} étale fini de degré }4\end{gathered}
\arrow[d,leftrightarrow] \& \\
\begin{gathered}
\deleted{(\underline J(\Sigma_{\underline J}),t\in\Gamma(\Sigma_J^*/S),\underline I\ \text{\deleted{rev.} torseur sous}}\\
(\underline I,\ \text{d'où }\underline J,\ t\in\Gamma(\Sigma_J^{\unsure{\alpha|*}}/S))
\end{gathered}\arrow[d,leftrightarrow]\arrow[r,leftrightarrow]\&
\begin{gathered}
\text{NB }\underline t\ \text{et }\underline I\ \text{indépendants}\\
(\underline J,\ (\text{d'où }\underline V_{\underline J}\ \text{et }\Sigma_J),\
\text{section }\inserted{t}\ \text{de }\Sigma_J^\alpha\ \text{sur }S\\
\text{torseur }I\ \text{sur }V_{\underline J})
\end{gathered}\\
\begin{gathered}
(\underline J,\ t\in\Gamma(\Sigma_{\underline J}^\alpha/S),\
X\ \text{rev. de }\Sigma_J^{\unsure{*}},\ \deleted{\unsure{g}}\ \text{étale sur }\Sigma_J^\alpha,\\
\text{jusqu'à isom sur }S,\ \text{sur chaque fibre géom. rev. \uline{abélien}}\\
\deleted{\text{de degré }4}\ \text{\inserted{\unsure{annulé|universel}, \unsure{par} \(2\)}}\ \text{\inserted{\unsure{fini|finie}, \unsure{univ.}\ldots}})\\
\text{\uline{NB} }X\ \text{et }t\ \text{sont indép. l'un de l'autre}
\end{gathered}\arrow[d,leftrightarrow] \& \\
\begin{gathered}
(\underline J,\ t\in\Gamma(\Sigma_J^\alpha/S),\ \deleted{\unsure{\underline I}}\ \underline K\ \text{torseur sous }\underline J)\\
\text{\uline{NB} }\underline K\ \text{et }t\ \text{sont indép l'un de l'autre.}
\end{gathered}\&
\end{tikzcd}}}

% margin note written sideways
\margin{\raggedright\scriptsize\(\deleted{f_i\ f_j\ f_k}\qquad x_i\ x_j\ x_k\)\\
\(V_{\underline J}(M_0)\)\\\(\unsure{V_{\underline J}(M)}\longrightarrow
\deleted{\unsure{4V_{\underline J}|\mathfrak S}}(M_0)\)\\\(\simeq\underline{\mathrm{End}}(M_0)\)\\
\(\hfill\rotatebox{90}{\(\in\)}\hspace{0.8em}\)\\
\(\hfill 1\hspace{0.8em}\)}

\hrulefill

Quand \(\underline I\) est muni d'une section, de sorte que \(\underline I\simeq\underline J\amalg S\), alors on a un iso. can \(X\simeq\Sigma_J\), d'où \(f_{\deleted{\unsure{I}}}:\Sigma_J\unsure{-|\longrightarrow}\Sigma_J\)
